\appendixchapter{Submanifolds}
\label{app:submanifolds}
\chaptermark{Submanifolds}
\section*{Introduction}
The inverse, implicit, and constant-rank theorems describe how nonlinear
equations can be straightened locally. Chapter~13 supplies the intrinsic
language in which that straightening becomes geometry:
\[
\begin{gathered}
\text{inverse/implicit/constant-rank theory}+\text{smooth manifolds}\\
\longrightarrow\text{submanifold geometry}.
\end{gathered}
\]
This appendix joins those two parts of the book. It distinguishes a smooth
parametrization from a well-behaved subset, identifies the tangent directions
allowed by constraints, and prepares the matrix groups of
Appendix~\ref{app:lie}. It is a continuation of Chapters~11 and~13, not a
prerequisite for Stokes.

Throughout this appendix manifolds are smooth, Hausdorff, second countable,
and without boundary; maps are smooth unless stated otherwise. Boundary
versions require additional hypotheses and are not implicit in our statements.
The primary comparisons are Tu~\cite[\S\S9,11]{tu-manifolds} and
Lee~\cite[Chapters 4--6]{lee-smooth}. We use the local slice viewpoint for
embedded subsets and reduce rank statements to Chapter~11, keeping the
topological requirement in an embedding separate from the rank requirement.

\section{Embedded Submanifolds as Local Straightening}
\label{sec:embedded-submanifolds}
A curved subset need not be a linear subspace. It can nevertheless have
exactly $k$ independent local coordinates if an ambient coordinate change
makes it a coordinate plane.
\begin{definition}[Embedded submanifold]\label{def:embedded-submanifold}
A subset $S\subset M^n$, with the subspace topology, is an
\textbf{embedded $k$-submanifold} if every $p\in S$ has a chart
$\varphi:U\to\varphi(U)\subseteq\R^n$ such that
\[
\varphi(U\cap S)=\varphi(U)\cap(\R^k\times\{0\}).
\]
Such a chart is called adapted to $S$. The induced coordinates on $S$ are
the first $k$ components of $\varphi|_{U\cap S}$.
\end{definition}
\begin{proposition}[The induced smooth structure]
These induced coordinates make $S$ a smooth $k$-manifold, and the inclusion
$i:S\to M$ is smooth with injective differential.
\end{proposition}
\begin{proof}
Identify the slice with $\R^k$. Its intersection with the open set
$\varphi(U)$ is open in that slice, so the induced coordinates are charts
for the subspace topology. An overlap map is obtained by inserting zeros,
applying an ambient smooth transition, and retaining the first $k$
components. It and its inverse are smooth. Hausdorffness and second
countability pass to subspaces: separating neighborhoods restrict to
the subspace (\cref{def:hausdorff}), and a countable basis restricts
as in \cref{prop:second-countable-inheritance}.
In these charts inclusion is $x\mapsto(x,0)$; its derivative
$v\mapsto(v,0)$ is injective. The induced maximal atlas is independent of
the chosen adapted charts, because the same overlap argument compares
any two of them.
\end{proof}
We henceforth identify $T_pS$ with the linear subspace $di_p(T_pS)$ of
$T_pM$. The identification is through this injective map, not an equality
of the original chart-equivalence-class constructions.

\begin{example}[Planes, open subsets, and graphs]
A coordinate plane is already straight. Every open subset of $M$ is an
embedded submanifold of dimension $\dim M$. For smooth
$f:U\subseteq\R^k\to\R^q$ with $U$ open, the graph
$\Gamma_f=\{(x,f(x)):x\in U\}$ is straightened by
\[
\varphi:U\times\R^q\to U\times\R^q,\qquad
\varphi(x,y)=(x,y-f(x)).
\]
Its inverse is $(x,z)\mapsto(x,z+f(x))$, so it is a diffeomorphism.
Graph coordinates give
\[
T_{(x,f(x))}\Gamma_f=\{(v,Df(x)v):v\in\R^k\}.
\]
Here $D$ is the Euclidean derivative; the displayed vector is the image
under the intrinsic differential of the graph parametrization. The
implicit function theorem, \cref{thm:implicit-function}, produces just
such graphs when an appropriate block of derivatives is invertible.
\end{example}
\begin{center}
\begin{tikzcd}[column sep=large,row sep=large]
\Gamma_f\cap(U\times\R^q)\arrow[r,hook]\arrow[d,"\varphi|"'] &
U\times\R^q\arrow[d,"\varphi"]\\
U\times\{0\}\arrow[r,hook] & U\times\R^q
\end{tikzcd}
\end{center}
The diagram records what straightening preserves: points of the subset
become exactly the zero-normal-coordinate points, while nearby ambient
points still have all $k+q$ coordinates.

\begin{proposition}[Testing a map with values in a submanifold]
For an embedded $S\subset M$, a map $h:P\to S$ is smooth if and only if
$i\circ h:P\to M$ is smooth.
\end{proposition}
\begin{proof}
One implication is composition. For the other, continuity into the
subspace $S$ follows from continuity into $M$. Near each image point use
an adapted chart. The coordinate expression of $h$ consists of the first
$k$ components of that of $i\circ h$, and is therefore smooth.
\end{proof}
This elementary test will verify the group operations in Appendix~\ref{app:lie}
without reconstructing their smoothness in every submanifold chart.

\section{Rank, Immersions, Submersions, and Embeddings}
\label{sec:manifold-rank}
For $f:M^m\to N^n$, the intrinsic differential
$df_p:T_pM\to T_{f(p)}N$ has the chart representative
$D(\psi\circ f\circ\varphi^{-1})(\varphi(p))$.
The coordinate conversions are isomorphisms, so rank is intrinsic.
\begin{definition}\label{def:manifold-immersion-embedding}
The map $f$ is an \textbf{immersion at $p$} if $df_p$ is injective,
and a \textbf{submersion at $p$} if it is surjective. It is an immersion
or submersion if the condition holds at every point.
A \textbf{smooth embedding} is an immersion that is a homeomorphism
from $M$ onto $f(M)$ with its subspace topology in $N$.
\end{definition}
An immersion preserves independent directions locally. An embedding also
requires the global subset topology to recover the topology of the source.

\begin{theorem}[Manifold constant-rank theorem]\label{thm:manifold-constant-rank}
If $f:M^m\to N^n$ has constant rank $r$ on a neighborhood of $p$, there
are charts centered at $p$ and $f(p)$ in which its expression is
\[
(x^1,\ldots,x^m)\longmapsto(x^1,\ldots,x^r,0,\ldots,0).
\]
\end{theorem}
\begin{proof}
The proof has two steps: express the intrinsic rank in coordinates, then
absorb Euclidean coordinate changes into the charts. Choose charts
$\varphi,\psi$ and shrink the source so that its image is in the target
chart. The Euclidean map $h=\psi\circ f\circ\varphi^{-1}$ has rank $r$
throughout this neighborhood, since its derivative represents $df$.
Apply \cref{thm:constant-rank} to obtain local Euclidean diffeomorphisms
straightening $h$. Composing them with $\varphi$ and $\psi$ gives the
required manifold charts. No new analytic assertion is needed.
\end{proof}
\begin{corollary}[The three maximal-rank models]
Near a point where $df_p$ is invertible, $f$ is a local diffeomorphism.
Near a submersion point it is a coordinate projection, and near an immersion
point it is the coordinate inclusion $x\mapsto(x,0)$.
\end{corollary}
\begin{proof}
A nonzero maximal minor remains nonzero in a neighborhood by continuity.
Rank there is therefore maximal and constant. Apply the theorem. In the
equal-dimensional case the coordinate model is the identity; the other
two cases are projection and inclusion.
\end{proof}
\begin{proposition}[Embeddings and embedded images]\label{prop:embedding-image}
A smooth embedding has an embedded image, and is a diffeomorphism onto
that image. Conversely the inclusion of an embedded submanifold is a
smooth embedding. In particular an injective immersion from a compact
manifold into a Hausdorff manifold is an embedding.
\end{proposition}
\begin{proof}
The issue is to exclude other pieces of the image from a local slice.
The immersion normal form makes $f(U)$ a slice near $f(p)$ for a source
neighborhood $U$. Since $f$ is a homeomorphism onto its image, $f(U)$ is
relatively open in $f(M)$. Intersect the target chart with an ambient open
set whose intersection with $f(M)$ lies in $f(U)$, and shrink again to a
coordinate box. The whole image in this box is now exactly its coordinate
slice. The inverse in slice coordinates is a coordinate projection and
is smooth. The converse follows from the induced-chart proposition.
Finally a continuous bijection from compact $M$ to its image in a Hausdorff
space is a homeomorphism by Chapter~10's compact-to-Hausdorff criterion;
combine this with the assumed immersion.
\end{proof}

\begin{example}[An immersed curve with a crossing]\label{ex:immersed-crossing}
The map $\gamma:S^1\to\R^2$ given in angle coordinates by
$\gamma(t)=(\sin t,\sin2t)$ is well defined and smooth.
Its derivative is $(\cos t,2\cos2t)$, never zero: when $\cos t=0$,
$\cos2t=-1$. Thus it is an immersion. The distinct circle points
$t=0,\pi$ both map to the origin, so it is not an embedding.
The two tangent lines there have slopes $2$ and $-2$. The image itself
is not an embedded curve: a small punctured neighborhood of the crossing
has four arms, whereas a punctured coordinate interval has two.
\end{example}
\begin{example}[Injectivity alone still does not suffice]
Restrict $\gamma(t)=(\sin t,\sin2t)$ to $(-\pi,\pi)$.
If two inputs have the same sine, either they agree or their cosines are
opposite. Equality of $\sin2t=2\sin t\cos t$ then forces either the
inputs to agree or $\sin t=0$; in this interval the latter gives $t=0$.
Thus the restriction is an injective immersion. But $t_j\to\pi$ from
below gives $\gamma(t_j)\to\gamma(0)$ without $t_j\to0$. Its inverse
on the image is not continuous, so it is not an embedding.
\end{example}

\begin{remark}[Immersion, injective immersion, embedding]
\[
\begin{array}{c|l}
\text{immersion}&df_p\text{ is injective at every point}\\
\text{injective immersion}&\text{also distinct points have distinct images}\\
\text{embedding}&\text{also the inverse onto the image is continuous}
\end{array}
\]
An immersion is a local differential condition. Injectivity removes
crossings, but an embedding also requires the source topology to agree
with the subspace topology on the image. The crossing and accumulation
examples above separate these requirements. For compact sources and
Hausdorff targets, \cref{prop:embedding-image} supplies the last condition.
\end{remark}

\section{Regular Values and Tangent Kernels}
\label{sec:manifold-regular-values}
Equations select subsets; their differentials select the allowed directions.
The rank hypothesis ensures that these two selections fit together smoothly.
\begin{definition}
A \textbf{regular point} of $f:M\to N$ is a point where $df_p$ is
surjective. A \textbf{regular value} $c\in N$ is a value for which every
point of $f^{-1}(c)$ is regular. An empty inverse image makes this condition
vacuously true. Points that are not regular are called critical points.
\end{definition}
\begin{theorem}[Regular level sets]\label{thm:manifold-regular-level}
Let $c$ be a regular value of $f:M^n\to N^m$. If $f^{-1}(c)$ is nonempty,
it is an embedded $(n-m)$-submanifold $S$, and, under the inclusion,
\[
\boxed{T_pS=\ker df_p\qquad(p\in S).}
\]
An empty level set is an embedded empty submanifold; no assertion of a
negative dimension is intended when $n<m$.
\end{theorem}
\begin{proof}
Near $p\in S$, surjectivity persists by a nonzero $m$-rowed minor.
The constant-rank theorem makes $f$ the projection onto its first $m$
coordinates, with $c$ represented by zero. Its fiber is the slice
$x^1=\cdots=x^m=0$; permuting coordinates puts this in the form of
\cref{def:embedded-submanifold}. Differentiating $f\circ i=c$ gives
$df_p\circ di_p=0$, hence $T_pS\subseteq\ker df_p$. Both spaces have
dimension $n-m$, the latter by rank--nullity, so they are equal.
\end{proof}
For Euclidean $f$, the identified kernel is $\ker Df(p)$, exactly the
formula of \cref{cor:regular-level-set}. The intrinsic theorem transports
that computation across charts rather than changing it.

\begin{example}[Spheres and cylinders]
For $h:\R^n\to\R$, $h(x)=\|x\|^2$, we have $Dh(p)v=2p\cdot v$.
At $h(p)=1$ this functional is nonzero. Consequently
\[
S^{n-1}=h^{-1}(1),\qquad T_pS^{n-1}=\{v:p\cdot v=0\}.
\]
Similarly $h(x,y,z)=x^2+y^2$ gives the cylinder $h^{-1}(1)$ and tangent
condition $xv_1+yv_2=0$, leaving the vertical direction free.
A rank failure is a failure of the criterion, not a proof of singularity:
the zero set of $x\mapsto x^2$ is still a point.
\end{example}
\begin{example}[An orthogonality constraint]\label{ex:orthogonal-level}
Let $M_n(\R)$ be the $n^2$-dimensional matrix space and let
$\operatorname{Sym}_n$ be its subspace of symmetric matrices, of dimension
$n(n+1)/2$. For $h:M_n(\R)\to\operatorname{Sym}_n$, $h(A)=A^TA$,
expanding $(A+tH)^T(A+tH)$ gives
\[
Dh(A)H=A^TH+H^TA.
\]
If $A^TA=I$, every symmetric $K$ is the image of $H=AK/2$.
Thus $I$ is a regular value with this specified target, and
\[
O(n)=h^{-1}(I),\quad \dim O(n)=\frac{n(n-1)}2,\quad
T_AO(n)=\{H:A^TH+H^TA=0\}.
\]
Taking all matrices as the target would incorrectly demand surjectivity
onto $n^2$ coordinates. Symmetry removes the redundant equations.
\end{example}

\begin{remark}[Three local pictures of the same submanifold]
\label{rem:submanifold-three-pictures}
The preceding results fit together as
\[
\text{coordinate slice}\ \longleftrightarrow\ \text{local graph}
\ \longleftrightarrow\ \text{regular level set}.
\]
For an embedded $k$-submanifold $S\subseteq M^n$, choose an adapted chart
$\varphi:U\to\varphi(U)\subseteq\R^n$. Projection onto the last $n-k$
coordinates gives a submersion $h=\operatorname{pr}_{\mathrm{last}}\circ\varphi$,
with
\[
S\cap U=h^{-1}(0),\qquad T_pS=\ker dh_p\quad(p\in S\cap U).
\]
Conversely, \cref{thm:manifold-regular-level} supplies an adapted slice for
a submersion's level set. The implicit function theorem, after selecting
an invertible derivative block in coordinates, expresses that level set
as a graph. A graph is straightened by $(x,y)\mapsto(x,y-f(x))$.
These are local descriptions in suitable coordinates, not a claim that an
entire submanifold is a single graph or a single global level set.
\end{remark}

\section{Constrained Extrema and a Global Perspective}
\begin{proposition}[Intrinsic multipliers]\label{prop:intrinsic-multipliers}
If $f:M\to\R$ has a local extremum on an embedded $S$ at $p$, then
$df_p$ vanishes on $T_pS$. If near $p$ the constraint is
$S=h^{-1}(c)$ for a submersion $h:M\to\R^m$, then there are scalars
$\lambda_1,\ldots,\lambda_m$ such that
\[
df_p=\sum_{j=1}^m\lambda_j\,dh^j_p.
\]
These multipliers are unique for this independent constraint list.
\end{proposition}
\begin{proof}
In coordinates on $S$, the ordinary necessary derivative condition for a
local extremum gives $d(f|_S)_p=0$. The chain rule says this is $df_p$
restricted to the image of $di_p$.
For the multiplier statement put $L=dh_p$. The kernel formula shows that
$df_p$ vanishes on $\ker L$. Since $L$ is onto, define a functional
$\ell:\R^m\to\R$ by $\ell(Lv)=df_p(v)$. Equal values of $Lv$ differ
by a kernel vector, so this definition is independent of the choice of
$v$ and is linear. Thus $df_p=\ell L$. Expanding $\ell$ in the standard
dual basis gives the displayed formula. Surjectivity of $L$ proves
uniqueness. This is quotient factorization in its simplest linear form.
\end{proof}
For $f(x)=a\cdot x$ on the unit sphere, the equation is $a=2\lambda p$.
For $a\ne0$ it gives $p=\pm a/\|a\|$, the maximum and minimum found by
Cauchy--Schwarz. The first-order equation supplies candidates; it does not
by itself decide which candidate is a maximum or a minimum.

\begin{theorem}[Whitney embedding theorem; stated without proof]
Every smooth $n$-manifold without boundary, for $n\geq1$, admits a smooth
embedding into $\R^{2n}$.
\end{theorem}
See Lee~\cite[Theorem 6.19]{lee-smooth}; Tu~\cite[\S11.3]{tu-manifolds}
discusses the embedding perspective. The theorem connects abstract and
embedded manifolds globally. Our local slice and rank proofs do not prove
it. Sard's theorem, transversality, tubular neighborhoods, Frobenius theory,
and Morse theory belong to subsequent study, signposted in Section~13.17.

\section*{Exercises}
\addcontentsline{toc}{section}{Exercises}
\begin{hexercise}{app-f-1}
Straighten the graph of $f(x,y)=x^2-y^2$ explicitly and compute its tangent
space at $(a,b,a^2-b^2)$. Compare the graph and regular-level calculations.
\end{hexercise}
\begin{hexercise}{app-f-2}
For $h(x,y)=xy$, determine all regular values and describe their level sets.
Explain precisely why the zero fiber fails to be an embedded curve at the origin.
\end{hexercise}
\begin{hexercise}{app-f-3}
Check the ranks of the projection $S^2\to\R$, $(x,y,z)\mapsto z$.
Compute the tangent line of a nonempty regular fiber as an intersection of kernels.
\end{hexercise}
\begin{hexercise}{app-f-4}
Verify the injectivity claim for the restricted crossing curve above and
exhibit the sequence witnessing failure of continuity of its inverse.
Explain why compactness prevents this phenomenon for an injective immersion.
\end{hexercise}
\begin{hexercise}{app-f-5}
For $1\leq k\leq n$, show that $h(A)=A^TA$ from real $n$ by $k$
matrices to symmetric $k$ by $k$ matrices has $I_k$ as a regular value.
Find the dimension and tangent space of the space of orthonormal $k$-frames.
\end{hexercise}
\begin{exercise}
Use intrinsic multipliers to find candidates for extrema of
$f(x,y,z)=xyz$ on the sphere. Explain why a necessary condition alone
does not classify all these candidates.
\end{exercise}
\begin{hexercise}{app-f-7}
Prove that the product of two embedded submanifolds is embedded by using
adapted product charts. Identify its tangent space and apply the result
to the formal torus $S^1\times S^1$.
\end{hexercise}
